% !TeX program = xelatex
\documentclass[11pt]{article}

% Self-contained, publisher-neutral Chinese cover letter.
\usepackage[margin=1in]{geometry}
\usepackage{fontspec}
\usepackage{xeCJK}
\usepackage{microtype}
\usepackage{xcolor}
\usepackage[hidelinks]{hyperref}
\input{%%AUTHOR_METADATA_PATH%%}

\definecolor{ink}{HTML}{20262E}
\definecolor{placeholder}{HTML}{6B7280}

\color{ink}
\setlength{\parindent}{0pt}
\setlength{\parskip}{0.72\baselineskip}
\pagestyle{empty}

% 灰色文字为写作提示；正式投稿前必须替换或删除。
\newcommand{\guidance}[1]{\textcolor{placeholder}{【#1】}}

\begin{document}

\hfill\today\par
\EditorName\par
\EditorTitle\par
\textit{\JournalName}

尊敬的编辑：

现随信提交题为《\ManuscriptTitle》的稿件，拟以\ArticleType 的形式投稿至
\textit{\JournalName}，恳请贵刊审议。

\guidance{首先用一至两句话介绍研究的科学背景或实际应用背景，并说明该问题
的重要性。} 尽管该领域近年来已取得一定进展，\guidance{指出本稿件所针对的
具体科学问题、现有研究局限或尚未解决的知识空白}。

在本研究中，我们\guidance{简要说明研究方法、研究设计、数据来源或理论框架，
避免加入过多技术细节}。研究结果表明，\guidance{用清晰、简洁的语言概括两至
三项最重要的研究发现}。这些发现\guidance{说明本研究在理论、方法或实际应用
方面的主要创新，以及该创新对相关领域的重要意义}。

我们认为，本研究与\textit{\JournalName}的办刊范围高度契合，因为
\guidance{说明稿件如何符合期刊的研究范围、重点方向或读者需求}。本研究结果
预计将受到\guidance{指出可能对此项研究感兴趣的研究领域、学术群体或专业人员}
的关注。

本稿件为原创成果，未曾在其他期刊或出版物上发表，目前也未同时投稿至其他
期刊。全体作者均已审阅并批准本稿件，同意将其投稿至\textit{\JournalName}。
\guidance{说明利益冲突；如无利益冲突，可写：“作者声明不存在任何利益冲突。”}
\guidance{如期刊要求，可补充作者贡献、伦理审批、数据可用性、推荐审稿人或
回避审稿人等声明。}

感谢您审阅本稿件。期待您的回复。

此致\par
敬礼！\par
\vspace{0.22in}

\begingroup
\setlength{\parskip}{2pt}
{\bfseries\CorrespondenceAuthorNameZh}\par
通讯作者，谨代表全体作者\par
\CorrespondenceAuthorAffiliationZh\par
\CorrespondenceAuthorAddress\par
电子邮箱：\href{mailto:\CorrespondenceAuthorEmail}{\CorrespondenceAuthorEmail}
\endgroup

\end{document}
